\documentclass[12pt,a4paper]{article}

% --- Lingua e font ---
\usepackage[utf8]{inputenc}
\usepackage[T1]{fontenc}
\usepackage[italian]{babel}
\usepackage{lmodern}
\usepackage{microtype}

% Impaginazione
\usepackage[a4paper,margin=2.5cm]{geometry}
\usepackage{parskip}
\linespread{1.15}
\usepackage{enumitem}
\setlist[itemize]{itemsep=3pt, topsep=5pt, leftmargin=*}

% Colori e titoli
\usepackage[dvipsnames]{xcolor}
\definecolor{primary}{RGB}{30,60,114}
\usepackage{titlesec}
\titleformat{\section}
{\Large\bfseries\color{primary}}{\thesection}{0.8em}{}[\vspace{-0.4em}\color{primary}\titlerule]
\titleformat{\subsection}
{\large\bfseries\color{primary}}{\thesubsection}{0.8em}{}
\titlespacing*{\section}{0pt}{1.6em}{0.8em}
\titlespacing*{\subsection}{0pt}{1.2em}{0.5em}

% Intestazione e piè di pagina
\usepackage{fancyhdr}
\setlength{\headheight}{15pt}
\pagestyle{fancy}
\fancyhf{}
\fancyhead[L]{\small\color{gray} RevHire -- Technical Report}
\fancyhead[R]{\small\color{gray}\leftmark}
\fancyfoot[C]{\small\thepage}
\renewcommand{\headrulewidth}{0.4pt}

% Collegamenti
\usepackage{hyperref}
\usepackage{color}
\hypersetup{
    colorlinks=true,
    linkcolor=primary,
    urlcolor=primary,
    pdftitle={RevHire -- Technical Report},
    pdfauthor={Alessandro Bonelli, Alessia Gatto, Elisa Picilli}
}

\begin{document}

%Copertina
    \begin{titlepage}
        \centering
        \vspace*{3cm}
        {\small\scshape Technical Report\par}
        \vspace{1cm}
        {\fontsize{40}{44}\selectfont\bfseries\color{primary} RevHire\par}
        \vspace{0.4cm}
        {\Large Piattaforma web per la ricerca di lavoro\par}
        \vspace{0.8cm}
        {\color{primary}\rule{0.5\textwidth}{0.8pt}\par}
        \vspace{0.8cm}
        {\large Corso di Sustainable Software Engineering\par}
        \vfill
        {\large
        Alessandro Bonelli [NF22500420]\\[0.3em]
        Alessia Gatto [NF22500351]\\[0.3em]
        Elisa Picilli [NF22500332]\par}
        \vspace*{3cm}
    \end{titlepage}

    \tableofcontents
    \newpage

% Contenuto
    \section{Introduzione}

    \subsection{Descrizione del progetto originale}
    RevHire è un'applicazione web dedicata alla ricerca di lavoro, realizzata in
    Java con il framework Spring Boot. Il progetto è stato sviluppato da un
    gruppo di studenti durante la seconda fase di un percorso di formazione
    (Revature Training Program).

    Gli utenti si dividono in due ruoli. Il candidato può registrarsi, compilare
    un curriculum (obiettivo, formazione, esperienze, competenze, progetti,
    certificazioni) oppure caricarlo in PDF o DOCX fino a 2~MB, cercare offerte
    con vari filtri e candidarsi con un click. Può anche ritirare una candidatura,
    salvare le offerte tra i preferiti e controllare lo stato di ciò che ha
    inviato.

    Il datore di lavoro, invece, registra la propria azienda, pubblica le
    offerte e le può modificare, chiudere, riaprire o eliminare. Per ogni
    offerta vede i candidati, può filtrarli per competenze ed esperienza,
    accettare o rifiutare le candidature e scrivere note interne. Ha anche una
    dashboard con il numero di offerte, di offerte attive, di candidature
    ricevute e di quelle ancora da valutare.

    L'applicazione segue il pattern MVC. Il codice è diviso in controller,
    servizi, repository ed entità, le pagine sono generate lato server con
    Thymeleaf e i dati sono salvati su un database Oracle tramite JPA/Hibernate.

    \subsection{Ambito e funzionalità del sistema}
    RevHire fa parte delle piattaforme di \emph{recruiting}: serve a far
    incontrare chi offre lavoro e chi lo cerca, e a gestire la candidatura dalla
    pubblicazione dell'offerta fino alla valutazione.

    Le funzioni principali sono:
    \begin{itemize}
        \item registrazione, accesso e gestione del profilo;
        \item gestione del curriculum;
        \item pubblicazione e gestione delle offerte;
        \item ricerca delle offerte con filtri;
        \item candidatura e controllo dello stato;
        \item valutazione dei candidati da parte del datore di lavoro;
        \item notifiche all'interno dell'applicazione.
    \end{itemize}

    \subsection{Tecnologie utilizzate}
    Il progetto è costruito con Maven, tramite il Maven Wrapper già incluso nel
    repository, e gira su Java con Spring Boot 3.2.5. Le tecnologie principali
    sono queste:
    \begin{itemize}
        \item \textbf{Spring Boot e Spring MVC}: struttura dell'applicazione e
        gestione delle richieste web;
        \item \textbf{Thymeleaf}: generazione delle pagine HTML lato server;
        \item \textbf{Spring Data JPA e Hibernate}: accesso ai dati e
        mappatura tra classi Java e tabelle;
        \item \textbf{Oracle Database}: database relazionale in cui sono salvati
        utenti, offerte e candidature;
        \item \textbf{Spring Security (modulo di cifratura)}: protezione delle
        password degli utenti;
        \item \textbf{OpenHTMLtoPDF e Apache PDFBox}: generazione di documenti
        in formato PDF;
        \item \textbf{Lombok}: riduzione del codice ripetitivo nelle classi Java.
    \end{itemize}

    Il progetto prevede inoltre dei test automatici, scritti con JUnit e
    Mockito, e il codice dei test è organizzato come quello dell'applicazione
    (controller, servizi e repository).

\end{document}